% 第9章：硬类别丢失概率信息。单个输入，正类概率 p 取 0.6 或 0.9。
% 源任务零一损失：p>0.5 时最优类别为 1（两种情形都存为“类别 1”）。
% 目标代价：标签0判1代价 0.8，标签1判0代价 0.2；R_T(1)=0.8(1-p)，R_T(0)=0.2p，阈值 0.8。
% p=0.6：R_T(1)=0.32, R_T(0)=0.12，应判 0；p=0.9：R_T(1)=0.08, R_T(0)=0.18，应判 1。
\begin{tikzpicture}[figure picture,x=1mm,y=1mm]
\def\L{75}
% 源任务决策带
\fill[FigureBlueFill] (0,21) rectangle (0.5*\L,25);
\fill[FigureRedFill] (0.5*\L,21) rectangle (\L,25);
\node[figure note,anchor=east] at (-2,23) {源任务};
\node[font=\figurefont\footnotesize] at (0.25*\L,23) {判 0};
\node[font=\figurefont\footnotesize] at (0.75*\L,23) {判 1};
% 目标任务决策带
\fill[FigureBlueFill] (0,3) rectangle (0.8*\L,7);
\fill[FigureRedFill] (0.8*\L,3) rectangle (\L,7);
\node[figure note,anchor=east] at (-2,5) {目标任务};
\node[font=\figurefont\footnotesize] at (0.4*\L,5) {判 0};
\node[font=\figurefont\footnotesize] at (0.9*\L,5) {判 1};
% 数轴
\draw[FigureStroke,line width=.75pt] (0,13) -- (\L,13);
\foreach \v in {0,0.5,0.8,1}{\draw[FigureStroke,line width=.6pt] (\v*\L,12) -- (\v*\L,14);}
\node[figure note,anchor=east] at (-2,13) {$p$};
\node[font=\figurefont\footnotesize,anchor=north] at (0,12) {0};
\node[font=\figurefont\footnotesize,anchor=north] at (\L,12) {1};
\draw[FigureStroke,line width=.6pt,dash pattern=on 2pt off 1.5pt] (0.5*\L,12) -- (0.5*\L,26);
\draw[FigureStroke,line width=.6pt,dash pattern=on 2pt off 1.5pt] (0.8*\L,1) -- (0.8*\L,14);
\node[figure note] at (0.5*\L,28) {源阈值 0.5};
\node[figure note] at (0.8*\L,-1.5) {新阈值 0.8};
% 两个可能的概率
\foreach \v/\t in {0.6/{$p=0.6$},0.9/{$p=0.9$}}{%
  \fill[FigureInk] (\v*\L,13) circle (1mm);
  \node[font=\figurefont\footnotesize] at (\v*\L,16.8) {\t};}
% 保存的硬类别
\node[figure note,align=left,anchor=west] at (\L+3,23) {两点都存为\\“类别 1”};
\node[figure note,align=left,anchor=west] at (\L+3,5) {两点要求\\相反决策};
\end{tikzpicture}
